\documentclass[12pt,a4paper]{article}

\usepackage[utf8]{inputenc}
\usepackage[T2A]{fontenc}
\usepackage[russian]{babel}
\usepackage{amsmath, amssymb, amsthm}
\usepackage{geometry}
\usepackage{algorithm}
\usepackage{algpseudocode}
\usepackage{graphicx}
\usepackage{hyperref}
\usepackage{booktabs}

\geometry{top=2cm, bottom=2cm, left=2.5cm, right=1.5cm}

\newtheorem{theorem}{Теорема}
\newtheorem{definition}{Определение}
\newtheorem{lemma}{Лемма}

\floatname{algorithm}{Алгоритм}

\title{Билет №80 \\ \small Информационно-поисковые системы. Классификация. Методы реализации и ускорения поиска. (ИТМО) \\ Информационно-поисковые системы. Классификация. Методы реализации и ускорения поиска. Методы построения и способы организации полнотекстовых индексов. (СпБГУ)}
\author{Сериков Владислав Александрович}
\date{16 мая 2026}

\begin{document}

\maketitle
\tableofcontents
\newpage

\section{Введение. Классификация информационно-поисковых систем}

\begin{definition}
\textbf{Информационно-поисковая система (ИПС)} --- это программно-аппаратный комплекс, предназначенный для хранения, поиска и выдачи неструктурированной или слабоструктурированной информации (обычно текстовой) в ответ на информационный запрос пользователя.
\end{definition}

\subsection{Классификация ИПС}

ИПС классифицируются по нескольким ключевым признакам:

\begin{enumerate}
    \item \textbf{По объекту поиска:}
    \begin{itemize}
        \item \textit{Документальные} --- результатом поиска является список документов (текстов, статей, веб-страниц), релевантных запросу.
        \item \textit{Фактографические} --- система извлекает конкретные факты из массивов текстов (например, "В каком году родился Ньютон?").
        \item \textit{Семантические} --- используют онтологии и графы знаний для понимания смысла запроса, а не просто совпадения слов.
    \end{itemize}

    \item \textbf{По масштабу и архитектуре:}
    \begin{itemize}
        \item \textit{Персональные} --- поиск по локальному диску (например, Spotlight в macOS).
        \item \textit{Корпоративные (Enterprise)} --- поиск по внутренним документам компании (Intranet).
        \item \textit{Глобальные (Web-scale)} --- веб-поисковики (Google, Яндекс), требующие распределенной архитектуры и краулинга.
    \end{itemize}

    \item \textbf{По модели поиска (Методы реализации):}
    \begin{itemize}
        \item Булева модель (Boolean Model).
        \item Векторная модель (Vector Space Model, VSM).
        \item Вероятностная модель (Probabilistic Model, BM25).
    \end{itemize}
\end{enumerate}

\section{Методы реализации поиска: Базовые модели}

Для реализации поиска необходимо математически описать соответствие между запросом $Q$ и документом $D$.

\subsection{Булева модель}
Документ рассматривается как множество терминов. Запрос формулируется с помощью логических операторов (\texttt{AND}, \texttt{OR}, \texttt{NOT}).
\textbf{Недостатки:} Отсутствие ранжирования (документ либо релевантен, либо нет), сложность формулирования точных запросов.

\subsection{Векторная модель (Vector Space Model)}
Каждый документ $D_j$ и запрос $Q$ представляются как векторы в $n$-мерном пространстве терминов, где $n$ --- размер словаря коллекции.
Вес термина $t_i$ в документе $D_j$ часто вычисляется по схеме \textbf{TF-IDF}:
$$w_{i,j} = \text{TF}(t_i, D_j) \times \text{IDF}(t_i) = f_{i,j} \times \log\left(\frac{N}{df_i}\right)$$
где $f_{i,j}$ --- частота термина в документе, $N$ --- общее число документов, $df_i$ --- количество документов, содержащих термин $t_i$.

Мера близости оценивается косинусом угла между векторами:
$$\text{sim}(Q, D_j) = \cos(Q, D_j) = \frac{\sum_{i=1}^{n} w_{i,q} w_{i,j}}{\sqrt{\sum_{i=1}^{n} w_{i,q}^2} \sqrt{\sum_{i=1}^{n} w_{i,j}^2}}$$

\section{Методы построения и способы организации полнотекстовых индексов}

Для эффективного поиска по большим коллекциям используется структура данных, называемая \textbf{обратным индексом} (Inverted Index).

\subsection{Структура обратного индекса}
Обратный индекс состоит из двух основных компонентов:
\begin{enumerate}
    \item \textbf{Словарь (Dictionary):} хранит все уникальные термины коллекции. Для каждого термина хранится указатель на его постинг-список и статистика (например, $df$). Часто реализуется в виде B-дерева, префиксного дерева (Trie) или хеш-таблицы.
    \item \textbf{Постинг-списки (Postings Lists):} для каждого термина хранится список идентификаторов документов (\textit{docID}), в которых он встречается. Постинг-списки всегда отсортированы по \textit{docID} для быстрого пересечения.
\end{enumerate}

\subsection{Алгоритм пересечения постинг-списков}
При обработке запроса "Термин1 \texttt{AND} Термин2" необходимо пересечь их постинг-списки.

\begin{algorithm}[H]
\caption{Пересечение двух отсортированных постинг-списков}
\begin{algorithmic}[1]
\Function{Intersect}{$p_1, p_2$}
    \State $answer \gets \text{пустой список}$
    \While{$p_1 \neq \text{NIL}$ \textbf{and} $p_2 \neq \text{NIL}$}
        \If{$docID(p_1) = docID(p_2)$}
            \State \Call{Add}{$answer, docID(p_1)$}
            \State $p_1 \gets \text{next}(p_1)$
            \State $p_2 \gets \text{next}(p_2)$
        \ElsIf{$docID(p_1) < docID(p_2)$}
            \State $p_1 \gets \text{next}(p_1)$
        \Else
            \State $p_2 \gets \text{next}(p_2)$
        \EndIf
    \EndWhile
    \State \Return $answer$
\EndFunction
\end{algorithmic}
\end{algorithm}

\textbf{Пример выполнения алгоритма (Трассировка):}\\
Пусть даны два постинг-списка:\\
$p_1 = [2, 4, 8, 16, 32]$\\
$p_2 = [1, 2, 3, 8, 9]$

\begin{itemize}
    \item \textit{Шаг 1:} Указатели на 2 ($p_1$) и 1 ($p_2$). $2 > 1 \implies$ сдвигаем $p_2$.
    \item \textit{Шаг 2:} Указатели на 2 ($p_1$) и 2 ($p_2$). Равны! Добавляем 2 в $answer$. Сдвигаем оба.
    \item \textit{Шаг 3:} Указатели на 4 ($p_1$) и 3 ($p_2$). $4 > 3 \implies$ сдвигаем $p_2$.
    \item \textit{Шаг 4:} Указатели на 4 ($p_1$) и 8 ($p_2$). $4 < 8 \implies$ сдвигаем $p_1$.
    \item \textit{Шаг 5:} Указатели на 8 ($p_1$) и 8 ($p_2$). Равны! Добавляем 8 в $answer$. Сдвигаем оба.
    \item \textit{Шаг 6:} Указатели на 16 ($p_1$) и 9 ($p_2$). $16 > 9 \implies$ сдвигаем $p_2$. Конец списка $p_2$.
    \item \textit{Результат:} $answer = [2, 8]$.
\end{itemize}

\subsection{Методы построения индекса: алгоритм SPIMI}
Для коллекций, не помещающихся в оперативную память, используется алгоритм \textbf{SPIMI} (Single-Pass In-Memory Indexing).

\textbf{Шаги алгоритма SPIMI:}
\begin{enumerate}
    \item Читать токены из потока документов, пока оперативная память не заполнится.
    \item Если термин встречается впервые, добавить его в словарь в памяти и создать новый постинг-список.
    \item Если термин уже есть, добавить текущий \textit{docID} в конец его постинг-списка (списки растут динамически, сортировка словаря на лету не требуется).
    \item Когда память заполнена, отсортировать словарь по алфавиту.
    \item Записать получившийся блок (индекс части коллекции) на диск. Очистить память.
    \item После обработки всей коллекции выполнить многопутевое слияние (multi-way merge) всех записанных на диск блоков в один финальный обратный индекс.
\end{enumerate}

\textbf{Минимальный пример шагов 1-5 (построение блока):}\\
Документ 1: "to do" $\implies$ токены: $(to, 1), (do, 1)$.\\
Документ 2: "to be" $\implies$ токены: $(to, 2), (be, 2)$.

\textit{Состояние памяти по мере чтения потока:}
\begin{itemize}
    \item Читаем $(to, 1)$: Словарь: $\{to: [1]\}$
    \item Читаем $(do, 1)$: Словарь: $\{to: [1], do: [1]\}$
    \item Читаем $(to, 2)$: Термин $to$ найден. Словарь: $\{to: [1, 2], do: [1]\}$
    \item Читаем $(be, 2)$: Словарь: $\{to: [1, 2], do: [1], be: [2]\}$
\end{itemize}
\textit{Сортировка и запись на диск:}
$be \rightarrow [2]$,
$do \rightarrow [1]$,
$to \rightarrow [1, 2]$.

\section{Методы ускорения поиска и сжатие}

Для ускорения обработки запросов на уровне структур данных и архитектуры применяются следующие методы.

\subsection{Skip-pointers (Указатели пропуска)}
В длинных постинг-списках последовательный проход неэффективен. Решение --- добавить указатели, позволяющие "перепрыгивать" через участки списка. Если текущий \textit{docID} значительно меньше искомого, система переходит по skip-указателю, избегая лишних сравнений. Оптимальный шаг между указателями --- $\sqrt{L}$, где $L$ --- длина списка.

\subsection{Эшелонирование (Tiered Indexes) и Чемпионские списки}
Вместо того чтобы искать по всему постинг-списку, список разделяется на уровни (эшелоны):
\begin{itemize}
    \item \textit{Уровень 1 (Champion list):} содержит $K$ документов с самыми высокими весами (например, TF-IDF) для данного термина.
    \item \textit{Уровень 2:} остальные документы.
\end{itemize}
Поиск сначала выполняется по первому уровню. Если набрано достаточное количество результатов, второй уровень не обрабатывается (Pruning).

\subsection{Сжатие индексов. Теория префиксного кодирования}
Постинг-списки хранят последовательности монотонно возрастающих целых чисел (docID). Вместо самих docID хранят дельты (разности) между соседними:
$[2, 5, 12, 14] \rightarrow [2, 3, 7, 2]$. Дельты --- это небольшие числа, которые можно сжать, используя коды переменной длины (Elias Gamma, Variable Byte).

Для того чтобы коды переменной длины однозначно декодировались без разделителей, они должны быть \textbf{префиксными} (ни одно кодовое слово не является началом другого). Фундаментальным ограничением на длины кодовых слов префиксного кода является теорема Крафта-Макмиллана.

\begin{theorem}[Неравенство Крафта --- Макмиллана]
Пусть задан алфавит из $D$ символов. Префиксный код с длинами кодовых слов $l_1, l_2, \dots, l_n$ существует тогда и только тогда, когда выполняется неравенство:
$$ \sum_{i=1}^{n} D^{-l_i} \le 1 $$
\end{theorem}

\begin{proof}
Докажем прямое утверждение для бинарного алфавита ($D=2$).
Сопоставим префиксному коду полное бинарное дерево. Корнем дерева является пустая строка. Переход к левому потомку означает добавление $0$, к правому --- $1$. Узлы дерева соответствуют префиксам кодовых слов.

Пусть максимальная длина кодового слова равна $l_{max}$. Построим полное бинарное дерево высоты $l_{max}$. Оно содержит ровно $2^{l_{max}}$ листьев.
Каждое кодовое слово длины $l_i$ соответствует некоторому узлу на глубине $l_i$. Поскольку код префиксный, ни одно кодовое слово не является предком другого кодового слова. Это значит, что если узел на глубине $l_i$ является кодовым словом, то все листья в поддереве, растущем из этого узла, не могут соответствовать никаким другим кодовым словам.

Поддерево с корнем на глубине $l_i$ содержит ровно $2^{l_{max} - l_i}$ листьев на максимальной глубине $l_{max}$.
Так как поддеревья для разных кодовых слов не пересекаются (из-за свойства префиксности), сумма листьев, покрытых всеми кодовыми словами, не может превышать общего количества листьев в дереве:
$$ \sum_{i=1}^{n} 2^{l_{max} - l_i} \le 2^{l_{max}} $$
Разделив обе части неравенства на $2^{l_{max}}$, получаем требуемое утверждение:
$$ \sum_{i=1}^{n} 2^{-l_i} \le 1 $$
\end{proof}

Основываясь на этой теореме, строятся оптимальные методы сжатия постинг-списков, такие как $\gamma$-код Элиаса (Elias Gamma), который кодирует число $x$ префиксом из унарной записи длины его бинарного представления и суффиксом из самого бинарного представления (без старшего бита).

\end{document}
